\documentclass[11pt]{article}
\usepackage{amsmath,amssymb,amsthm}
\usepackage[margin=1.1in]{geometry}
\usepackage{hyperref}

\newtheorem{theorem}{Theorem}
\newtheorem{lemma}{Lemma}
\newtheorem{proposition}{Proposition}
\newtheorem{remark}{Remark}
\theoremstyle{definition}
\newtheorem{definition}{Definition}

\newcommand{\T}{\mathbb T}
\newcommand{\R}{\mathbb R}
\newcommand{\Z}{\mathbb Z}
\newcommand{\Q}{\mathbb Q}
\newcommand{\PP}{\mathbb P}
\newcommand{\NN}{\mathcal N}
\newcommand{\Cov}{\mathrm{Cov}}

\title{Spectral-front identities for periodic Navier--Stokes\\
and an obstruction to uniform pointwise Osgood closure}
\author{ns-singularity-certificate-lab\thanks{Machine-assisted
research repository; all exact-arithmetic certificates cited below are
reproducible from the repository at the commit recorded in the
appendix. This draft claims no resolution of any Clay Millennium
problem, in either direction.}}
\date{August 2026 --- draft}

\begin{document}
\maketitle

\begin{abstract}
For strong solutions of the unforced incompressible Navier--Stokes
equations on $\T^3$ we prove that the spectral bandwidth
$N_0^2=\|\nabla u\|_2^2/\|u\|_2^2$ obeys the differential inequality
$\frac{d}{dt}\log N_0^2\le K D/(2\nu)$, where
$D=\|\nabla u\|_2^2$ and
$K=\|\PP(u\cdot\nabla u)\|_2^2/\|\nabla u\|_2^4$, with an explicit
nonnegative defect; consequently if
$\int_0^{T_{\max}}KD\,dt<\infty$ the solution is global, and in
particular any finite-time breakdown forces the action to diverge. The
action admits the exact representation
$\int_0^{T'}KD\,dt=\int_0^{T'}\|\partial_tu\|_2^2/D\,dt
+\nu^2\int_0^{T'}\|{-}\Delta u\|_2^2/D\,dt+\nu\log(D(T')/D(0))$; by
AM--GM this yields the one-sided criterion
$\int KD\le2\int(\|\partial_tu\|_2^2/D+\nu^2\|\Delta u\|_2^2/D)$
(the converse comparison fails for explicit decaying solutions), and
$KD$ is dominated pointwise by both the Serrin $(\infty,2)$ and the
critical $\|\nabla u\|_{L^3}^2$ actions. On the obstruction side we construct an
explicit family of divergence-free trigonometric fields with exact
closed-form spectral laws, and prove for it the capacity lower bound
$\|\PP(u_N\cdot\nabla u_N)\|_2^2\ge c_0N^3$ with the sharp exponent,
uniformly in the cutoff profile and in the seed vector. It follows
\emph{unconditionally} that no Osgood-admissible pointwise field
inequality $K\le\Phi(\log N_0^2)$ can hold uniformly: the dichotomy
cannot be upgraded to unconditional regularity through that route. The
capacity constant $c_0$ and the band threshold $N_*$ are the only
non-effective constants in the paper; every other constant is explicit.
Proofs are elementary and self-contained modulo classical local
theory; no numerical statement enters any proof, and all numerics are
confined to a certificate appendix. The claimed
novelty is the exact packaging --- the modal identities, their
machine-checkable defect ledger, and the obstruction family with its
capacity bound --- and overlap with known refinements of the Serrin
criteria is possible. Internal review was performed by AI agents and is
not human peer review.
\end{abstract}

\section{Introduction}

The regularity theory of the three-dimensional Navier--Stokes
equations is governed by critical quantities whose a-priori control is
unknown. We isolate a scalar action, $\int K D\,dt$, that (i) bounds
the growth of the spectral bandwidth through an exact modal identity
with computable defects, (ii) is dominated by the classical critical
actions, so its finiteness is a weaker hypothesis than theirs, and
(iii) possesses an exact dynamic representation splitting it into a
supercritical part (the $\dot H^1$-bandwidth action) and a harmless
logarithm. We then show, \emph{unconditionally}, that
the hypothesis cannot be verified by any pointwise Osgood-type field
inequality: an explicit smoothly truncated coherent critical-spectrum
family obeys an exact capacity lower bound with the sharp exponent,
which forces any uniform majorant
$\Phi(\log N_0^2)$ of $K$ to be of exponential class. The pair of
results freezes this lane: the dichotomy is at least as strong as the
classical criteria it refines, and the pointwise-Osgood route to
making it unconditional is closed.

We emphasise what is \emph{not} claimed: no global regularity for all
data, no blow-up, no Clay statement. Numerics never enter any proof;
they appear only as reproducible exact-arithmetic certificates
corroborating the proven statements.

\section{Setting and main results}

Let $\T^3=(\R/2\pi\Z)^3$ with normalised measure,
$u(x,t)=\sum_{k\in\Z^3\setminus\{0\}}\hat u_k(t)e^{ik\cdot x}$ real,
zero-mean, divergence-free ($k\cdot\hat u_k=0$), solving
$\partial_tu+\PP(u\cdot\nabla u)=\nu\Delta u$ with $\nu>0$ and
$u_0\in H^m_\sigma$, $m>5/2$, $u_0\neq0$; let $u$ be the maximal
strong solution on $[0,T_{\max})$. Write
$H_r=\sum_k|k|^{2r}|\hat u_k|^2$, $D=H_1$, $N_0^2=H_1/H_0$,
$N_1^2=H_2/H_1$, $z=\log N_0^2\ (\ge0)$,
$\NN=-\PP(u\cdot\nabla u)$, and
$K=\|\NN\|_2^2/D^2$.

\begin{theorem}\label{main}
With the notation above:
\begin{itemize}
\item[(a)] $\displaystyle
\Lambda(t)=\log N_0^2(t)-\frac1{2\nu}\int_0^tKD\,ds$ is
non-increasing on $[0,T_{\max})$; quantitatively
$\frac{d}{dt}\log N_0^2=\frac{2\Cov}{N_0^2}-\frac{2\nu V}{N_0^2}
\le\frac{KD}{2\nu}$ with
$\Cov=T_1/H_0$, $V=\mathrm{Var}_p(|k|^2)$,
$p(k)=|\hat u_k|^2/H_0$, and an explicit nonnegative two-part defect.
\item[(b)] Exactly one of: $\int_0^{T_{\max}}KD\,dt<\infty$, and then
$z$ is bounded, $u\in L^\infty(0,T_{\max};H^1)$ and $T_{\max}=\infty$;
or $\int_0^{T_{\max}}KD\,dt=\infty$. In particular
$T_{\max}<\infty\Rightarrow\int_0^{T_{\max}}KD\,dt=\infty$; no claim
is made that a global solution has finite action.
\item[(c)] For every $T'<T_{\max}$,
\[
\int_0^{T'}KD\,dt=\int_0^{T'}\frac{\|\partial_tu\|_2^2}{D}\,dt
+\nu^2\int_0^{T'}N_1^2\,dt+\nu\log\frac{D(T')}{D(0)} ,
\]
and by AM--GM
$\int_0^{T'}KD\,dt\le2\int_0^{T'}(\|\partial_tu\|_2^2/D+\nu^2N_1^2)dt$;
the reverse comparison fails ($u=e^{-\nu t}(0,0,\cos x_1)$ has
$\int KD=0$ with divergent bandwidth action).
\item[(d)] Pointwise, $KD\le\|u\|_{L^\infty}^2$ and
$KD\le C_S^2\|\nabla u\|_{L^3}^2$.
\item[(e)] If moreover $KD\le\Phi(z)D+R$ a.e.\ with $\Phi>0$
nondecreasing, $\int^\infty ds/\Phi=\infty$, $R\ge0$,
$\int_0^{T_{\max}}R\,dt<\infty$, then the first alternative of (b)
holds.
\end{itemize}
\end{theorem}

The proof (Section~\ref{proofs}) is a chain of eight elementary lemmas;
the only external inputs are the classical local theory, the
subcritical Serrin class $L^\infty L^6$, and the mean-zero Sobolev
embedding.

On the obstruction side, we use a \emph{smoothly} truncated coherent
family.

\begin{definition}[admissible cutoff]\label{admissible}
$\chi\in C^\infty([0,\infty);[0,1])$ with $\chi\equiv1$ on
$[0,\tfrac12]$ and $\operatorname{supp}\chi\subset[0,1]$; no
monotonicity is assumed. Write $c_\chi=1+\|\chi'\|_\infty$.
\end{definition}

Fix $v_0\in\R^3\setminus\{0\}$ (no integrality is needed) and an
admissible $\chi$, and let
\begin{equation}\label{family}
\hat u_N(k)=\chi\!\Big(\frac{|k|}{N}\Big)\frac{P_kv_0}{|k|^2}\ \ (k\neq0),
\qquad \hat u_N(0)=0,\qquad
P_k=I-\frac{k\otimes k}{|k|^2},
\end{equation}
a finite trigonometric polynomial banded at $|k|\le N$.

\begin{lemma}[exact family laws]\label{familylaws}
$u_N$ is real, zero-mean and divergence-free, and with
$S_N^\chi=\sum_{k\neq0}\chi(|k|/N)^2|k|^{-2}$,
$T_N^\chi=\sum_{k\neq0}\chi(|k|/N)^2|k|^{-4}$,
$\Sigma_N^\chi=\sum_{k\neq0}\chi(|k|/N)|k|^{-2}$:
$H_0=\tfrac23\|v_0\|^2T_N^\chi$, $H_1=\tfrac23\|v_0\|^2S_N^\chi$,
$N_0^2=S_N^\chi/T_N^\chi$, $u_N(0)=\tfrac23\Sigma_N^\chi v_0$, and
$\|u_N\|_\infty^2/\|\nabla u_N\|_2^2\ge2(\Sigma_N^\chi)^2/(3S_N^\chi)$.
Moreover, for $N\ge32$ and every admissible $\chi$,
\[
\frac{N}{348}\le S_N^\chi\le128N,\qquad 6\le T_N^\chi\le128,\qquad
\frac{N}{44544}\le N_0^2\le\frac{64}{3}N ,
\]
all crude explicit constants, proven; so $N_0^2\asymp N$.
\end{lemma}

The band-symmetry identity behind Lemma~\ref{familylaws} is needed for a
\emph{general} radial weight $f$ with $\sum|k|^2|f|<\infty$ --- no
positivity and no sharp ball --- since $\chi$ is not assumed monotone.

\begin{theorem}[capacity lower bound; sharp exponent]\label{capacity}
For every $v_0\in\R^3\setminus\{0\}$ and every admissible $\chi$ there
exist $c_0=c_0(\chi,v_0)>0$ and $N_*=N_*(\chi,v_0)$ with
\[
\|\PP(u_N\cdot\nabla u_N)\|_2^2\;\ge\;c_0N^3\qquad\text{for all }N\ge N_* .
\]
The exponent $3$ is sharp, since (d) with Lemma~\ref{familylaws} gives
$\|\PP(u_N\cdot\nabla u_N)\|_2^2\le\|u_N\|_\infty^2H_1\asymp N^3$.
\end{theorem}

Theorem~\ref{capacity} replaces the capacity \emph{hypothesis} of earlier
versions of this paper. Its proof is given in full in the repository file
\texttt{lstar/lstar\_proof\_main.md} (Theorem 7.1(3)) and is summarised in
Section~\ref{proofs}; we state here only what it delivers.

\begin{theorem}[no Osgood-admissible uniform majorant]\label{nogo}
If $\Phi$ is nondecreasing and
$K(u)\le\Phi(\log N_0^2(u))$ for every zero-mean divergence-free real
trigonometric field $u$, then $\Phi(s)\ge ce^{s}$ for all large $s$;
in particular $\int^\infty ds/\Phi<\infty$, so clause (e) of
Theorem~\ref{main} cannot be activated by any uniform pointwise field
inequality.
\end{theorem}

\begin{remark}[what is open, and no longer used]\label{lstarremark}
The \emph{sharply} truncated family $\hat u_N(k)=P_kv_0/|k|^2$ on
$1\le|k|\le N$ satisfies the same exact laws with $S_N,T_N$ in place of
$S_N^\chi,T_N^\chi$ and saturates the Bernstein ratio two-sidedly,
$\tfrac23S_N\le\|u_N\|_\infty^2/\|\nabla u_N\|_2^2\le S_N$. The
corresponding capacity bound for \emph{that} family --- this paper's
former Hypothesis~(L\textsuperscript{*}) --- remains \textbf{open}. It is
now \textbf{unused}: Theorem~\ref{nogo} quantifies over all real
zero-mean divergence-free trigonometric fields, so the smooth family
alone suffices. The honest statement is ``the obstruction is
unconditional'', not ``(L\textsuperscript{*}) is proven''. A proven
negative result constrains the sharp family further: no constant-vector
sweeping split, in any H\"older pairing and any constant pairing
direction, can exceed $\|\PP(u_N\cdot\nabla u_N)\|_2^2\gtrsim N$ ---
two powers short --- which is why a concentrated test field, not a
constant sweep, is the right instrument.
\end{remark}

\begin{remark}[effectivity]\label{effectivity}
The constants $c_0$ and $N_*$ of Theorem~\ref{capacity}, and hence $c$
of Theorem~\ref{nogo}, are \textbf{non-effective}: the test field driving
the capacity bound is produced from the density of
$C^\infty_{c,\sigma}(\R^3)$ in $L^2_\sigma(\R^3)$, a pure existence
argument. Every other constant in this paper is explicit. An effective
version would require an explicitly constructed divergence-free $\Psi$
with $\int(V\cdot\nabla V)\cdot\Psi\neq0$ together with a quantitative
lower bound on that pairing; non-vanishing is certified only on
$\{0<|\zeta|<\pi^2/2048\}$ intersected with two cone conditions, so such
a $\Psi$ has spatial radius $R\gtrsim10^3$ and forces $N_*\gtrsim10^4$.
\end{remark}

\section{Proofs}\label{proofs}

The full lemma chain L0--L11 with complete proofs is reproduced from
the repository file \texttt{complete\_proof.md}; we give it here in
compressed form.

\paragraph{L0a (regularity of moments).}
$H_0,H_1,H_2$ are continuous; $H_0,H_1\in C^1$ with
$\dot H_0=2\langle u,\partial_tu\rangle$,
$\dot H_1=2\langle-\Delta u,\partial_tu\rangle$ (difference quotients
in $C^1H^{m-2}$ paired against $C(L^2)$);
$\partial_tu=\NN-\nu(-\Delta)u$ in $L^2$. Positivity of $H_0$ on all
of $[0,T_{\max})$ is \emph{not} elementary and is proven as L0b
\emph{after} L7, by bootstrap: if $t_1$ is the first vanishing time,
then on $[0,t_1)$ Lemmas 1--4 and 7 apply, $KD\le\|u\|_\infty^2\le
C_mM^2$ with $M=\sup_{[0,t_1]}\|u\|_{H^m}$, so $N_0^2$ is bounded
there, $\int_0^{t_1}N_0^2<\infty$, and the integrated energy identity
gives $H_0(t_1)>0$, a contradiction.

\paragraph{L1 (ledger; energy neutrality).}
$\tfrac12\dot H_r=T_r-\nu H_{r+1}$ ($r=0,1$) with
$T_r=\sum|k|^{2r}\mathrm{Re}\langle\hat u_k,\hat\NN_k\rangle$;
$T_0=-\langle u,(u\cdot\nabla)u\rangle=0$; hence
$\int_0^{T_{\max}}D\,dt\le H_0(0)/(2\nu)$.

\paragraph{L2 (identity).}
$\frac{d}{dt}\log N_0^2=2\Cov/\mu-2\nu V/\mu$ with $\mu=N_0^2$:
subtract the logarithmic derivatives and use
$H_0H_2-H_1^2=H_0^2V$.

\paragraph{L3 (closable covariance bound).}
$|\Cov|\le\sqrt{V\|\NN\|_2^2/H_0}$: modal Cauchy--Schwarz
$|\mathrm{Re}\langle\hat u_k,\hat\NN_k\rangle|\le
|\hat u_k||\hat\NN_k|$, then vector Cauchy--Schwarz.

\paragraph{L4 (square completion; (a)).}
$\sup_{V\ge0}[\alpha\sqrt V-\nu V]=\alpha^2/(4\nu)$ with
$\alpha=\sqrt{\|\NN\|^2/H_0}$ gives
$\frac{d}{dt}\log N_0^2\le\|\NN\|_2^2/(2\nu H_1)=KD/(2\nu)$; the
defect is $\Gamma^{\rm CS}+\Gamma^{\rm SC}\ge0$ as displayed in the
repository file, with rational total.

\paragraph{L5 (dichotomy; (b)).}
$\int KD<\infty\Rightarrow z\le z_*\Rightarrow
H_1\le e^{z_*}H_0(0)\Rightarrow u\in L^\infty L^6$ (mean-zero Sobolev),
a subcritical Serrin class $\Rightarrow$ regularity and extension.

\paragraph{L6 (action representation; (c)).}
Expand $\|\NN\|_2^2=\|\partial_tu+\nu(-\Delta)u\|_2^2$, use
$2\nu\langle\partial_tu,-\Delta u\rangle=\nu\dot D$, divide by $D>0$,
integrate; $\int\dot D/D=\log(D(T')/D(0))$.

\paragraph{L7 (dominations; (d)).}
$\|\NN\|_2\le\|u\|_\infty\|\nabla u\|_2$ and
$\|\NN\|_2\le\|u\|_{L^6}\|\nabla u\|_{L^3}
\le C_S\|\nabla u\|_2\|\nabla u\|_{L^3}$.

\paragraph{L8 (Osgood closure; (e)).}
$\Omega(z)=\int_0^zds/\Phi$;
$\frac{d}{dt}\Omega(z)\le D/(2\nu)+R/(2\nu\Phi(0))$, both integrable;
$z$ bounded; apply L5.

\paragraph{L9--L11 and Theorem~\ref{nogo}.}
Band symmetry $\sum_{k\neq0}k_ik_jf=\delta_{ij}\tfrac13\sum|k|^2f$ for
any radial $f$ with $\sum|k|^2|f|<\infty$ (sign flips kill $i\neq j$;
permutations equalise the diagonal; no positivity is used); the family
laws of Lemma~\ref{familylaws} follow by applying it with the weights
$\chi(|k|/N)^2|k|^{-6}$, $\chi(|k|/N)^2|k|^{-4}$ and
$\chi(|k|/N)|k|^{-4}$ to
$|P_kv_0|^2=\|v_0\|^2-(k\cdot v_0)^2/|k|^2$ and to
$u_N(0)=\sum\hat u_N(k)$. The lattice bounds follow from the dyadic
shell count $\#\{2^j\le|k|<2^{j+1}\}<64\cdot8^j$ (upper) and an
all-positive sub-box count in $N/2<|k|\le N$ (lower), using
$\chi\equiv1$ on $[0,\tfrac12]$ for $S_N^\chi\ge S_{\lfloor N/2\rfloor}$
and the six unit vectors for $T_N^\chi\ge6$; the two-sided bound gives
$2^q$-adic gaps $c_-\le s_{2^qN}-s_N\le c_+$ with $q=20$,
$c_-=\log(1048576/950272)=0.0985\ldots$,
$c_+=\log(2^{20}\cdot950272)=27.63\ldots$, elementary and $N$-uniform.
By Theorem~\ref{capacity} and $H_1\le\tfrac{256}{3}\|v_0\|^2N$,
$K(u_{N_j})\ge c_0''e^{s_{N_j}}$ along $N_j=2^{20j}N_\sharp$;
monotonicity of $\Phi$ interpolates across the bounded gaps, so
$\Phi(s)\ge c_0''e^{-c_+}e^{s}$ for all large $s$ and
$\int^\infty ds/\Phi<\infty$.

\paragraph{Theorem~\ref{capacity} (summary; full proof in the
repository).}
Because $P_k=P_{k/N}$, the family obeys the exact scaling
$\hat u_N(k)=N^{-2}F(k/N)$ with $F(\xi)=\chi(|\xi|)P_\xi v_0/|\xi|^2$,
so with $V=\widetilde F$ (inverse transform on $\R^3$) the lattice sums
defining $u_N(y/N)$ and $\nabla u_N(y/N)$ are \emph{literally} Riemann
sums of spacing $1/N$:
$u_N(y/N)=N(V(y)+E_N(y))$, $\nabla u_N(y/N)=N^2(\nabla V(y)+E'_N(y))$
identically, with $\sup_{|y|\le R}(|E_N|+|E'_N|)=O(N^{-1}\log N)$
explicitly (the integrand is compactly supported with a single
integrable algebraic singularity; the gradient half needs and receives
its own, log-free, argument). Poisson summation is \emph{not} available
here --- $V$ decays exactly like $|y|^{-1}$ with a sign-definite
spherical mean, so the periodisation diverges and, its terms being
eventually positive, no summation method rescues it. Pairing against the
concentrated divergence-free test field $\psi_N(x)=N^{3/2}\Psi(Nx)$,
whose $L^2(\T^3)$ norm is $N$-independent, and using self-adjointness of
$\PP$ gives
$\|\PP(u_N\cdot\nabla u_N)\|_2^2\ge(2\pi)^{-3}N^3|I_\Psi+\mathcal
E_N|^2/\|\Psi\|_{L^2(\R^3)}^2$ with
$I_\Psi=\int(V\cdot\nabla V)\cdot\Psi$ and $\mathcal E_N=O(\log N/N)$.
It remains to find one $\Psi$ with $I_\Psi\neq0$, i.e.\ to show
$\PP(V\cdot\nabla V)\not\equiv0$. This holds for \emph{every} admissible
$\chi$ and \emph{every} $v_0$, because the leading small-$\zeta$
singularity of $F*F$ is the $\chi$-independent Oseen convolution $\tau$,
for which, with $v_0=\|v_0\|e_3$,
\[
\zeta\times\big(\tau(\zeta)\zeta\big)
=\frac{3\pi^3}{8}\|v_0\|^2\frac{\zeta_3}{|\zeta|}(\zeta\times e_3)\neq0
\]
off the two degenerate directions --- the Fourier form of the elementary
fact that the Oseen field is not a stationary Euler flow, since
$\nabla\times(W\times\nabla\times W)=12y_3(e_3\times y)/|y|^6$. A crude
perturbation bound $|T-\tau|\le64\pi\|v_0\|^2$ on $|\zeta|\le\tfrac18$
transfers this to the cutoff family on $|\zeta|<\pi^2/2048$.
\qed

\section{Position among known criteria; limitations}

By (d), the finiteness of $\int KD$ is implied by the classical Serrin
$(\infty,2)$ action and by $\int\|\nabla u\|_{L^3}^2dt$; the criterion
(b) is therefore at least as strong as both (we do not claim strictly
stronger). By (c) and its AM--GM corollary, the criterion is implied
by the a-priori control of
$\int(\|\partial_tu\|_2^2/D+\nu^2N_1^2)\,dt$ --- the
half-derivative-supercritical $\dot H^1$-bandwidth action --- while
the exact identity shows the two differ only by $\nu\log D$, which can
diverge to $-\infty$ on decaying solutions; this makes precise both
why the criterion is strong and why its unconditional verification is
out of reach of the present methods.
Theorem~\ref{nogo} closes the pointwise-Osgood route unconditionally.
Its one limitation is effectivity (Remark~\ref{effectivity}): the
capacity constant $c_0$ and the threshold $N_*$ are not produced, only
shown to exist. All statements are elementary; the
contribution is the exact packaging, the obstruction family with its
capacity bound, and the machine-verifiable defect ledger accompanying
(a).

\appendix
\section{Certificate appendix (numerics, separated from all proofs)}

No number below enters any proof. Both Theorem~\ref{main} and
Theorem~\ref{nogo} are unconditional, so --- unlike in earlier versions
--- these values no longer carry a hypothesis; they are corroboration.

\paragraph{Smooth family (the family used by Theorem~\ref{nogo}).}
Reproducible from the repository
(\texttt{experiments/run\_smooth\_family\_capacity.py}; JSON with
SHA-256 sidecar at
\texttt{outputs/lstar/smooth\_family\_capacity.json}). Seed
$v_0=(1,2,3)$; weight
$\chi(r)=1-S(\mathrm{clamp}((r^2-\tfrac14)/\tfrac34,0,1))$ with $S$ the
degree-9 smoothstep $126s^5-420s^6+540s^7-315s^8+70s^9$, which is a
polynomial in $r^2$ --- hence $\chi(|k|/N)\in\Q$ at every lattice point
and the whole family is exactly rational --- and lies in $C^4$, not
$C^\infty$. Definition~\ref{admissible} asks for $C^\infty$; the proof
uses smoothness of $\chi$ only through $c_\chi=1+\|\chi'\|_\infty$, so
$C^1$ suffices there, but as the class is written the exact lanes below
are evidence for the $C^p$ variant. Degree-$(2p+1)$ smoothsteps give
rational $\chi\in C^p$ for every $p$.

Exact-\texttt{Fraction} verification of the family laws at
$N=4,6,8,10,12$: every residual of $H_0$, $H_1$ and $u_N(0)$ is the
rational $0$, with $u_N(0)$ exactly parallel to $v_0$. Exact capacity by
full $O(|B_N|^2)$ contraction over all ordered pairs of active band
points:
\[
\begin{array}{r|rrrrrr}
N & 4 & 6 & 8 & 10 & 12 & 16\\\hline
\|\PP(u_N\!\cdot\!\nabla u_N)\|_2^2/N^3
 & 678.28 & 1027.31 & 1241.51 & 1384.36 & 1485.78 & 1619.65\\
K & 0.6166 & 1.1359 & 1.6705 & 2.2077 & 2.7464 & 3.8254\\
K/N_0^2 & 0.2662 & 0.3287 & 0.3643 & 0.3861 & 0.4008 & 0.4195
\end{array}
\]
A dealiased FFT continuation (grid $4N+2$, whole product spectrum
represented) extends this to $N=48$:
$\|\PP(u_N\cdot\nabla u_N)\|_2^2/N^3$ rises $1241.5\to1911.3$ and
$K/N_0^2$ rises $0.3643\to0.4574$ over $N=8\to48$, both monotonically.
Exact versus FFT at the two shared bands agrees to double-precision
roundoff: relative difference $5.5\times10^{-16}$ ($N=8$) and
$7.0\times10^{-16}$ ($N=16$) in the norm, $5.3\times10^{-16}$ and
$8.1\times10^{-16}$ in $K$. Two-point Richardson diagnostics (not
proofs): $\|\PP(u_N\cdot\nabla u_N)\|_2^2/N^3\to2064.3$,
$K/N_0^2\to0.4764$, $H_1/N\to87.51$, $N_0^2/N\to0.5671$ --- positive
limits, i.e.\ the sharp exponent, consistent with
Theorem~\ref{capacity}. The continuum non-degeneracy
$\PP(V\cdot\nabla V)\not\equiv0$ is corroborated three independent ways
(exact rational lattice modes; float 3-D quadrature of the continuum
convolution; an axisymmetric radial reduction giving
$\mathrm{curl}(V\cdot\nabla V)=Z(r)\sin\theta\cos\theta\,e_\phi$ with
$Z\to-12\pi^4\|v_0\|^2/r^4$), agreeing to $10^{-4}$ or better.

\paragraph{Sharp family (corroboration only; family no longer used).}
These are \emph{sharp}-family numbers, supporting the open
Hypothesis~(L\textsuperscript{*}) of Remark~\ref{lstarremark}, which no
statement in this paper uses. Reproducible from
\texttt{experiments/run\_coherent\_family\_certificates.py} and
\texttt{run\_osgood\_gate.py} (JSON with SHA-256 sidecars under
\texttt{outputs/verification\_sprint\_v1/osgood\_gate/}):
exact-\texttt{Fraction} family laws at $N=4,6,8$ (residual $0$); exact
$K(u_4)=0.78841070\ldots$ (full fraction stored),
$K(u_6)=1.43447181\ldots$, $K(u_8)=2.03724301\ldots$; float pipeline
agreement $0.0$, $4.6\times10^{-16}$, $1.3\times10^{-15}$;
$Q(v_0)>0$ exact at $N=4,6,8$ with float continuation
$Q/N^2=850.7\to1029.5$ over $N=8\to32$; captured fraction of
$\|\PP(u\cdot\nabla u)\|_2$ by the sweeping pairing:
$0.921/0.917/0.915$ (exact, $N=4,6,8$) and
$0.915\to0.912$ (float, $N=8\to32$); $K/N_0^2=0.259\to0.396$
($N=4\to32$); Leray retention $0.780\to0.867$.

\end{document}
